\section{Educação financeira e matemática financeira no Ensino Médio}
\label{sec-fundamentacao-educacao-financeira}

A educação financeira abrange a compreensão de situações que envolvem o uso do dinheiro, o planejamento de gastos e a avaliação de decisões relacionadas ao consumo e ao crédito. A matemática financeira fornece instrumentos para analisar essas situações, como o cálculo de porcentagens e juros e a comparação de valores ao longo do tempo. O Banco Central do Brasil aborda a gestão do orçamento e o uso do crédito como temas da formação financeira \cite{bcb2026caderno}.


No Ensino Médio, essa relação encontra respaldo na Base Nacional
Comum Curricular. A habilidade EM13MAT203 contempla o uso de
conceitos matemáticos em atividades de planejamento e análise
com aplicativos e planilhas, incluindo situações de orçamento
e simulação de juros \cite{brasilBNCC2018}. Para o trabalho em
sala, essa orientação permite formular tarefas em que o estudante
escolhe entre alternativas, registra os cálculos realizados e
justifica sua decisão com base nos valores obtidos.


\section{Operações bancárias e movimentação de saldo}
\label{sec-fundamentacao-operacoes}

O saldo expressa a posição da conta em determinado instante.
O extrato registra as movimentações que explicam a passagem de
um saldo a outro. Se \(S_0\) representa o saldo inicial e \(m_k\)
representa o valor da \(k\)-ésima movimentação, com sinal positivo
para entradas e negativo para saídas, o saldo após \(n\)
movimentações é dado por
\[
    S_n = S_0 + \sum_{k=1}^{n} m_k.
\]

Por exemplo, uma conta com saldo inicial de \mbox{R\$~200,00}
que registra um pagamento de \mbox{R\$~45,00} e recebe uma
transferência de \mbox{R\$~30,00} passa a apresentar saldo de
\mbox{R\$~185,00}. A leitura do extrato permite reconstruir esse
resultado e explicar por que o saldo diminuiu em
\mbox{R\$~15,00}, embora a conta tenha recebido um crédito.

\section{Juros, prazo e comparação de valores}
\label{sec-fundamentacao-juros}



Em uma operação de crédito, o valor da prestação informa quanto
será pago em cada vencimento. Para examinar o custo da decisão,
é preciso considerar também o número de pagamentos, o valor
financiado e os encargos incluídos. Em um financiamento com
\(n\) prestações iguais de valor \(P\), sem entrada ou cobranças
adicionais, o total pago corresponde a \(nP\). Duas propostas
podem, portanto, apresentar prestações de valores próximos e
totais distintos quando seus prazos diferem. A metodologia da
Calculadora do Cidadão relaciona prestação, prazo, taxa mensal
e valor financiado em operações com capitalização mensal
\cite{bcbMetodologiaFixas}.

Nas aplicações financeiras, a questão de interesse é a evolução
do capital investido. Como referência matemática, um capital
inicial \(C\), aplicado a uma taxa composta constante \(i\) por
\(n\) períodos, sem aportes, retiradas ou encargos, resulta no
montante \(M=C(1+i)^n\). A unidade do período deve corresponder
à periodicidade da taxa. Uma taxa mensal, por exemplo, exige que
o prazo seja expresso em meses.

Esses conceitos orientam perguntas distintas em uma atividade
escolar. No crédito, o estudante pode investigar quanto pagará
ao longo do prazo escolhido. No investimento, pode examinar como
o valor aplicado evolui durante a atividade. No EDBank, a leitura
dos resultados de investimento deve considerar a modalidade
simulada e as regras de rendimento adotadas pelo sistema.


\section{Simulação digital como recurso didático}
\label{sec-fundamentacao-simulacao}

Em uma atividade de matemática financeira, a simulação permite
alterar uma condição da operação e examinar seus efeitos sobre o
resultado. Ao comparar dois cenários, o estudante precisa identificar
quais valores foram mantidos, qual parâmetro foi modificado e como
essa mudança explica a diferença observada.

Uma experiência brasileira com essa abordagem é apresentada por
Lima e Batista em um guia didático elaborado a partir de uma
intervenção com 33 estudantes do Ensino Médio de uma escola estadual
de Alagoas. A proposta articula o registro de despesas, a análise de
padrões de consumo e a simulação de aportes periódicos. Em uma das
etapas, os estudantes utilizam os dados produzidos pelo aplicativo
para fundamentar suas respostas. A interpretação dos resultados
integra, assim, a própria atividade de matemática financeira
\cite{limaBatista2026}.

No EDBank, uma atividade com finalidade semelhante pode partir da
compra de um mesmo produto em três ou quatro parcelas. O professor
pode pedir que os estudantes antecipem o total a pagar em cada
condição, consultem os valores apresentados pelo simulador e
expliquem a diferença a partir da regra de acréscimo adotada pelo
shopping. A tarefa relaciona porcentagem, número de parcelas e
custo total em uma situação que o estudante pode reproduzir.